% LTeX: enabled=false
\documentclass[11pt,a4paper,oneside]{book}
% --- Idioma y codificación ---
\usepackage[T1]{fontenc}
\usepackage[spanish]{babel}
\usepackage{titlesec}

% Elimina la etiqueta (número) del capítulo, pero conserva el título
\titleformat{\chapter}[block]
  {\normalfont\huge\bfseries} % formato del título
  {}                          % etiqueta vacía: aquí iría "Capítulo N"
  {0pt}                       % separación entre etiqueta y título
  {}                          % código previo al título

% --- Márgenes (ajusta a tu gusto) ---
\usepackage[a4paper,margin=2.5cm,headheight=14pt]{geometry}

% --- Código ---
\usepackage{listings}
\usepackage{xcolor}

\definecolor{codebg}{RGB}{248,248,248}
\definecolor{codekeyword}{RGB}{0,0,160}
\definecolor{codecomment}{RGB}{0,110,0}
\definecolor{codestring}{RGB}{150,0,0}

\lstset{
  backgroundcolor=\color{codebg},
  basicstyle=\ttfamily\small,       % tamaño del código
  keywordstyle=\color{codekeyword}\bfseries,
  commentstyle=\color{codecomment}\itshape,
  stringstyle=\color{codestring},
  numbers=left,                     % números de línea
  numberstyle=\tiny\color{gray},
  numbersep=8pt,
  frame=single,
  rulecolor=\color{gray!40},
  breaklines=true,                  % ¡clave para imprimir!
  showstringspaces=false,
  tabsize=2,
  captionpos=b,
  extendedchars=true,
  inputencoding=utf8,
}

% --- Entorno para cada programa ---
\newenvironment{programa}[2]{%
  \section{#1}%
  \noindent #2\par\medskip
}{%
  \medskip
}

% --- Navegación en PDF (opcional pero útil) ---
\usepackage[hidelinks]{hyperref}

\begin{document}

\frontmatter
\title{Notebook de BolitosFC}
\author{Camilo Sodo, Gaspar Hilman, Juan Ignacio Varisco}
\date{\the\year}
\maketitle
\tableofcontents

\mainmatter
% LTeX: enabled=true

\chapter{Generalidades}

\begin{programa}{Plantilla}{Plantilla de código para competencias.}
    \lstinputlisting[language=C++]{../template.cpp}
\end{programa}

\begin{programa}{Makefile}{Makefile para compilar el codigo mas comodamente.}
    \lstinputlisting[language=make]{../makefile}
\end{programa}

\begin{programa}{Busqueda Binaria}{Plantilla de una busqueda binaria basica.}
    \lstinputlisting[language=C++]{../binarySearch.cpp}
\end{programa}

\chapter{Grafos}

\begin{programa}{Recorridos}{Plantillas para recorridos BFS y DFS.}
    \lstinputlisting[language=C++]{../Grafos/template.cpp}
\end{programa}

\begin{programa}{Diametro en Arboles}{Calcula el camino mas largo de un árbol (su diametro).}
    \lstinputlisting[language=C++]{../Grafos/DiametroArboles.cpp}
\end{programa}

\begin{programa}{Dijkstra}{Busca para todos los vertices la distancia minima desde un punto de partida dado.}
    \lstinputlisting[language=C++]{../Grafos/Dijkstra.cpp}
\end{programa}

\begin{programa}{Ascenso Binario}{En arboles, permite encontrar un ancestro arbitrario de un vertice cualquiera rapidamente.}
    \lstinputlisting[language=C++]{../Grafos/BinaryLifting.cpp}
\end{programa}

\begin{programa}{Primer Ancestro Común}{En arboles, permite encontrar rapidamente el menor ancestro común (lowest common ancestor) entre dos vertices arbitrarios. Utiliza ascenso binario.}
    \lstinputlisting[language=C++]{../Grafos/PrimerAncestroComun.cpp}
\end{programa}

\begin{programa}{Componentes Furtemente Conexas}{Permite encontrar las componentes fuertemente conexas de un digrafo.}
    \lstinputlisting[language=C++]{../Grafos/CompFuertementeConexas.cpp}
\end{programa}

\begin{programa}{Ordenamiento Topologico}{Ordena los vertices de un DAG (digrafo aciclico dirigido) de forma que toda arista e = (v->w) cumple que v precede a w en el orden. Sirve para determinar si el grafo no es aciclico.}
    \lstinputlisting[language=C++]{../Grafos/OrdenTopologico.cpp}
\end{programa}

\chapter{Teoria de Números}

\begin{programa}{Minimo Comun Multiplo/Divisor Comun Menor}{Calcula el mcm o dcm de dos numeros.}
    \lstinputlisting[language=C++]{../Teoria de Numeros/mcm-dcm.cpp}
\end{programa}

\begin{programa}{Criba de Eratóstenes}{Calcula para todos los numeros menores a un N si son primos y su menor divisor.}
    \lstinputlisting[language=C++]{../Teoria de Numeros/CribaEratostenes.cpp}
\end{programa}

\begin{programa}{Exponenciación Binaria}{Calcula rapidamente la enecima potencia de un numero. Util para conseguir el inverso modular.}
    \lstinputlisting[language=C++]{../Teoria de Numeros/ExpBinaria.cpp}
\end{programa}

\chapter{Estructuras de Datos}

\begin{programa}{Segment Tree}{Permite realizar cualquier operacion asociativa en rangos.}
    \lstinputlisting[language=C++]{../Estructuras/SegmentTree.cpp}
\end{programa}

\begin{programa}{Lazy Segment Tree}{Permite actualizar rangos rapidamente.}
    \lstinputlisting[language=C++]{../Estructuras/LazySegmentTree.cpp}
\end{programa}

\begin{programa}{Trie}{Conjunto de strings con acceso e insercion rapida.}
    \lstinputlisting[language=C++]{../Estructuras/Trie.cpp}
\end{programa}

\begin{programa}{Union-Find}{Implementacion de conjuntos que optimizan la busqueda y la union.}
    \lstinputlisting[language=C++]{../Estructuras/UnionFind.cpp}
\end{programa}

\chapter{Códigos de Ejemplo}

\begin{programa}{Codigo 1}{}
    \lstinputlisting[language=C++]{../Ejemplos/cubrecadena.cpp}
\end{programa}

\begin{programa}{Codigo 2}{}
    \lstinputlisting[language=C++]{../Ejemplos/substring_problem_v3.cpp}
\end{programa}

\begin{programa}{Codigo 3}{}
    \lstinputlisting[language=C++]{../Ejemplos/substring_problem.cpp}
\end{programa}

\backmatter
\end{document}